\section{Introdução}

O Modelo de Atores de computação foi estabelecido por Carl Hewitt em 1973, baseado em trabalhos anteriores de Peter Bishop e Richard Steiger. Seus objetivos eram abordar os novos desafios impostos pelos sistemas massivamente concorrentes e paralelos que emergiram com a computação em nuvem e arquiteturas many-core.

Seus objetivos foram alcançados criando um modelo computacional que é inerentemente concorrente. Uma grande mudança dos modelos anteriores é que ele é baseado em leis físicas, ao invés de álgebra pura. A hipótese mais importante defendida pelo trabalho de Hewitt é:

\begin{quote}
Toda computação fisicamente possível pode ser diretamente implementada com Atores
\end{quote}

O Modelo de Atores representa uma reformulação completa da computação concorrente, passando de máquinas de estado globais para sistemas distribuídos, assíncronos e baseados em passagem de mensagens que melhor refletem a realidade física dos sistemas de computação modernos.
